%% 34장 방정식의 근
\chapter{Roots of Equations}
\label{ch:roots}

\abstract{$f(x)=0$을 푸는 일은 수치해석에서 가장 오래된 문제의 하나이다. 이 장에서는 근을 낀 구간을 반씩 줄이는 이분법의 확실함과, 도함수나 이전 점들을 써서 근으로 빠르게 다가가는 Newton 방법과 할선법의 속도를 견준다. 이어서 두 장점을 합친 Brent의 방법과 TOMS 748을 본다. 수렴 차수를 정의하고 Newton 방법의 2차 수렴, 할선법의 황금비 차수를 유도한다. 중근에서는 수렴이 1차로 떨어지는 이유와, 시작점이 나쁘면 Newton 방법이 두 점 사이를 끝없이 맴도는 예도 본다. 근의 조건수가 도달할 수 있는 정확도를 정한다는 것을 전개한 다항식으로 보이고, 마지막으로 방정식계로 넘어가 Newton 방법, 야코비안을 다시 계산하지 않는 Broyden의 갱신, MINPACK의 Powell 혼성법을 다룬다. scipy의 C 코드에는 컴파일러가 넣은 FMA가 있어서 stardust가 scipy의 근 찾기와 비트까지 같은 결과를 내려면 그 FMA를 명시적으로 적어야 했다. 이 일도 소개한다.}

\section{Bracketing}
\label{sec:bracket}

\subsection{Bisection}

$f$가 $[a,b]$에서 연속이고 $f(a)$와 $f(b)$의 부호가 다르면 중간값 정리에 의해 그 사이에 근이 있다. 가운데 점 $m=(a+b)/2$에서 $f$의 부호를 보고 근이 있는 반쪽을 남기면 구간의 길이가 반으로 준다. 이것이 \term{이분법}{bisection}이다. $k$걸음 뒤 구간의 길이는 $(b-a)2^{-k}$이므로, 허용오차 $\delta$에 이르는 걸음 수는 $\lceil\log_2((b-a)/\delta)\rceil$로 미리 정해진다. 함수가 아무리 사납게 굴어도 이 수는 변하지 않는다.

binary64에서는 이분법을 더 철저하게 할 수 있다. 같은 부호의 binary64 수는 비트 패턴을 정수로 읽은 순서가 값의 순서와 같으므로(8장) 가운데 점을 값의 평균 대신 정수의 평균으로 잡으면 구간 안의 binary64 수의 개수가 걸음마다 반으로 준다. 수는 많아야 $2^{64}$개이므로 $64$걸음 안에 이웃한 두 binary64 수까지 좁혀진다. 근을 하나의 ulp 안으로 확정해야 할 때, 예컨대 참값을 만들 때 쓰는 방법이다.

\subsection{Stopping Criteria}

실제 구현은 구간의 길이가 $\mathrm{xtol}+\mathrm{rtol}\cdot|x|$보다 작아지면 멈춘다. 절대 허용오차만 쓰면 근이 아주 클 때 binary64의 간격보다 작은 허용오차를 요구하게 되고, 상대 허용오차만 쓰면 근이 $0$일 때 끝나지 않는다. scipy의 구간법은 기본값이 $\mathrm{xtol}=2\times10^{-12}$, $\mathrm{rtol}=4\varepsilon$이다. 근이 $1$ 근처이면 $\mathrm{xtol}$이 지배하므로 결과는 대략 $2\times10^{-12}$의 정확도만 약속한다는 점에 주의해야 한다.

\section{Order of Convergence and Newton's Method}
\label{sec:newton}

\subsection{Order of Convergence}

근 $x_\ast$로 수렴하는 수열 $x_k$의 오차 $e_k=x_k-x_\ast$가
\[
\lim_{k\to\infty}\frac{|e_{k+1}|}{|e_k|^p}=C,\qquad0<C<\infty
\]
를 만족하면 수렴의 \term{차수}{order}가 $p$라 한다. $p=1$이고 $C<1$이면 1차 수렴이다. 이때 걸음마다 오차가 일정한 비율로 준다. 이분법은 오차의 상계가 걸음마다 반으로 주므로 1차와 같은 속도이다. $p=2$이면 맞는 자릿수가 걸음마다 대략 두 배가 된다.

\subsection{Newton's Method}

$x_k$에서 $f$를 접선으로 근사하고 그 접선의 근을 다음 점으로 삼는 것이 \term{Newton 방법}{Newton's method}이다.
\[
x_{k+1}=x_k-\frac{f(x_k)}{f'(x_k)}
\]

\begin{theorem}
\label{thm:newton-quadratic}
$f$가 근 $x_\ast$ 근처에서 두 번 연속 미분 가능하고 $f'(x_\ast)\ne0$, $f''(x_\ast)\ne0$이면, $x_0$가 $x_\ast$에 충분히 가까울 때 Newton 방법은 2차 수렴하고
\[
\lim_{k\to\infty}\frac{e_{k+1}}{e_k^2}=\frac{f''(x_\ast)}{2f'(x_\ast)}
\]
이다.
\end{theorem}

\begin{proof}
Taylor 정리에 의해 $0=f(x_\ast)=f(x_k)-f'(x_k)e_k+f''(\xi_k)e_k^2/2$인 $\xi_k$가 $x_k$와 $x_\ast$ 사이에 있다. 양변을 $f'(x_k)$로 나누고 정리하면
\[
e_{k+1}=e_k-\frac{f(x_k)}{f'(x_k)}=\frac{f''(\xi_k)}{2f'(x_k)}e_k^2
\]
이다. $x_\ast$를 중심으로 하는 작은 구간에서 $|f''/(2f')|\le M$이고 $M|e_0|<1$이면 $|e_{k+1}|\le(M|e_k|)|e_k|$이므로 오차가 단조롭게 줄어 수렴하고, 극한은 $\xi_k,x_k\to x_\ast$에서 나온다.
\end{proof}

$f(x)=\cos x-x$를 $x_0=1$에서 풀면 오차가 $1.1\times10^{-2}$, $2.8\times10^{-5}$, $1.7\times10^{-10}$이고, 네 번째 걸음에서 binary64로 정확한 근에 닿는다. 맞는 자릿수가 두 배씩 느는 것이 보인다.

\begin{aside}[Newton과 Raphson]
Newton 방법이라는 이름에도 불구하고 Newton의 원래 방법은 지금과 모습이 많이 달랐다. Newton은 다항식 $x^3-2x-5$를 풀면서 근의 근삿값에 보정을 더한 새 다항식을 만들고 그 다항식의 선형 부분만 풀었다. 도함수라는 개념은 쓰지 않았다. Raphson이 1690년에 매번 원래 다항식으로 돌아가는 꼴로 다듬었다. 오늘날의 $f/f'$ 꼴은 Simpson이 1740년에 일반 함수로 넓힌 것이다. Ypma가 이 과정을 자세히 정리했다\cite{ypma1995}.
\end{aside}

\subsection{Failures of Newton's Method}

\cref{thm:newton-quadratic}은 국소적인 정리이다. 시작점이 근에서 멀면 Newton 방법은 엉뚱한 곳으로 가거나 맴돈다. 고전적인 예는 $f(x)=x^3-2x+2$이다. $x_0=0$에서 시작하면 $x_1=1$, $x_2=0$, $x_3=1$로 두 점 사이를 영원히 오간다. 실수 근은 $-1.77$ 근처에 하나 있는데 그곳으로는 결코 가지 않는다. 또 $f'(x_k)$가 $0$에 가까우면 걸음이 터무니없이 커진다. 그래서 실용적인 근 찾기는 Newton 방법을 그대로 쓰지 않는다. 근을 낀 구간을 유지하며 걸음이 구간을 벗어나면 이분법으로 돌아가는 식으로 보호한다.

근이 중근이면 $f'(x_\ast)=0$이어서 정리의 가정이 깨진다. 중복도가 $m$인 근에서는 $f(x)\approx c(x-x_\ast)^m$이므로 $f/f'\approx(x-x_\ast)/m$이고, $e_{k+1}\approx(1-1/m)e_k$로 1차 수렴한다. $f(x)=(x-1)^3e^x$를 $x_0=2$에서 풀면 이웃한 오차의 비가 $0.733$, $0.718$, $0.705$, $0.695$, $0.687$, $0.681$로 $2/3$에 다가간다. 중복도를 알면 걸음에 $m$을 곱해 2차 수렴을 되찾을 수 있다.

\subsection{Halley's Method}

2차 도함수까지 쓰면 3차 수렴을 얻는다. 접선 대신 $f$와 두 번 접하는 쌍곡선의 근을 다음 점으로 삼는 \term{Halley 방법}{Halley's method}
\[
x_{k+1}=x_k-\frac{2f(x_k)f'(x_k)}{2f'(x_k)^2-f(x_k)f''(x_k)}
\]
이 그것이다\cite{scavo1995halley}. 도함수를 계산하는 비용이 함수 자체와 비슷하다면 걸음마다 일이 늘어나 Newton 방법보다 반드시 낫지는 않다. 그러나 특수함수의 역함수를 구할 때처럼 $f'$과 $f''$이 $f$의 계산에서 거의 공짜로 나오는 경우에는 Halley 방법이 강력하다. 39--41장의 특수함수의 역함수가 그런 예이다.

\section{Derivative-Free Methods}
\label{sec:derivative-free}

\subsection{The Secant Method}

도함수 대신 직전 두 점을 잇는 직선, 곧 할선의 근을 다음 점으로 삼으면 \term{할선법}{secant method}
\[
x_{k+1}=x_k-f(x_k)\frac{x_k-x_{k-1}}{f(x_k)-f(x_{k-1})}
\]
을 얻는다. 걸음마다 함수를 한 번만 부르므로 Newton 방법보다 싸다. 오차의 점화식은 29장의 나눗셈 차분으로 쓸 수 있다. 할선은 $x_{k-1}$, $x_k$의 1차 보간 다항식이므로 그 근에서 생기는 보간 오차를 쓰면
\[
e_{k+1}=\frac{f[x_{k-1},x_k,x_\ast]}{f[x_{k-1},x_k]}e_ke_{k-1}\approx\frac{f''(x_\ast)}{2f'(x_\ast)}e_ke_{k-1}
\]
이다. $|e_{k+1}|\approx A|e_k|^p$로 두고 대입하면 $A|e_k|^p\approx C|e_k|\cdot(|e_k|/A)^{1/p}$이므로 지수를 비교해 $p=1+1/p$, 곧 $p=(1+\sqrt5)/2\approx1.618$이다. 할선법의 수렴 차수는 황금비이다. 같은 $\cos x-x$를 $x_0=1$, $x_1=0.5$에서 풀면 오차가 $1.4\times10^{-2}$, $8.2\times10^{-4}$, $2.5\times10^{-6}$, $4.5\times10^{-10}$, $2.2\times10^{-16}$이다. 함수 호출 한 번당 얻는 자릿수로 견주면 할선법($1.618$)이 Newton 방법($\sqrt2\approx1.414$, 걸음마다 $f$와 $f'$을 부르므로)보다 효율적이다. 이 비교의 지표를 \term{효율 지수}{efficiency index}라 한다. Ostrowski가 도입했고 Traub가 반복법의 분류에 체계적으로 썼다\cite{traub1964}.

\subsection{Inverse Interpolation}

세 점을 쓰면 2차 보간을 할 수 있는데, $x$의 함수로 $f$를 보간하면 근을 구하려고 2차 방정식을 풀어야 한다. 대신 점 $(f_i,x_i)$를 지나는 2차식 $x=q(y)$를 만들어 $f$의 함수로 $x$를 보간하면 근의 근사는 $q(0)$으로 바로 나온다. 이것이 \term{역 2차 보간}{inverse quadratic interpolation}이다. 29장의 Lagrange 꼴로 쓰면
\[
q(0)=\frac{x_0f_1f_2}{(f_0-f_1)(f_0-f_2)}+\frac{x_1f_0f_2}{(f_1-f_0)(f_1-f_2)}+\frac{x_2f_0f_1}{(f_2-f_0)(f_2-f_1)}
\]
이다. 수렴 차수는 약 $1.839$이다.

\subsection{Brent's Method}

할선법과 역보간은 빠르지만 근을 낀 구간을 지키지 않아 발산할 수 있고, 이분법은 확실하지만 느리다. Dekker는 둘을 섞어 근을 낀 구간을 유지하면서 할선 걸음이 구간 안에 떨어지면 받아들이고 아니면 이분하는 방법을 만들었다\cite{dekker1969}. 그러나 Dekker의 방법은 할선 걸음이 아주 조금씩만 구간을 줄이는 경우에 이분법보다 훨씬 느려질 수 있다. Brent는 할선 대신 가능하면 역 2차 보간을 썼다. 또 보간 걸음이 직전 걸음들보다 충분히 줄지 않으면 이분법으로 바꾸는 조건을 더해 함수 호출의 수가 이분법 걸음 수의 제곱을 넘지 않도록 보장했다\cite{brent1973}. 실제로는 거의 언제나 초선형으로 수렴한다. $\cos x-x$의 근을 $[0,1]$에서 $10^{-15}$까지 구하는 데 이분법은 함수를 $52$번 부르지만 Brent의 방법은 $9$번 부른다.

Alefeld, Potra, Shi의 TOMS 748은 역 3차 보간과 이차 Newton 걸음을 섞고 구간이 반으로 줄지 않으면 이분해 근을 낀 구간을 유지하면서 효율 지수를 할선법보다 높였다\cite{alefeld1995}. 같은 문제에서 함수를 $8$번 부른다. scipy의 배열 원소별 근 찾기는 또 다른 혼합인 Chandrupatla의 방법을 쓴다\cite{chandrupatla1997}.

\section{Attainable Accuracy}
\label{sec:root-accuracy}

근을 낀 구간을 쓰는 방법이 쓰는 정보는 $f$의 부호뿐이고 다른 방법도 계산한 $f$ 값에 기대므로, 계산한 $f$ 값의 부호를 믿을 수 없는 구간 안에서는 어떤 방법도 근을 더 정확히 찾을 수 없다. $f$를 계산한 값에 크기 $\delta$의 오차가 있으면 단순근 근처에서 그 구간의 폭은 대략 $\delta/|f'(x_\ast)|$이다. 따라서 근의 조건수는 $1/|f'(x_\ast)|$이다. 중복도가 $m$인 근에서는 $f\approx c(x-x_\ast)^m$이므로 그 폭이 $(\delta/|c|)^{1/m}$으로 커진다.

$(x-1)^7$을 전개한 다항식 $x^7-7x^6+21x^5-35x^4+35x^3-21x^2+7x-1$을 Horner 방법으로 계산해 보자. 계산값의 오차는 대략 $10^{-14}$ 크기이므로 근 근처의 폭은 $(10^{-14})^{1/7}\approx0.01$이다. 실제로 $[0.98,1.02]$의 고른 점 $2001$개에서 계산값의 부호가 $438$번 바뀌고 그 점들이 $0.991$에서 $1.008$까지 퍼져 있다. scipy의 Brent 방법은 이 다항식의 근을 $0.99438$로 낸다. 오차 $5.6\times10^{-3}$은 방법이 아니라 문제의 탓이다. 다항식의 근이 계수의 작은 변화에 얼마나 민감한지는 Wilkinson이 체계적으로 분석했는데, 그의 유명한 예는 근이 $1,2,\dots,20$인 다항식으로 단순근조차 계수의 반올림에 크게 흔들린다는 것을 보였다\cite{wilkinson1959polynomials}. 같은 근을 $(x-1)^7$의 꼴로 계산하면 아무 문제가 없다. 근을 구하기 전에 함수를 어떻게 계산할지부터 따져야 한다.

\section{Systems of Equations}
\label{sec:systems}

\subsection{Newton's Method}

$F:\mathbb R^n\to\mathbb R^n$의 근을 찾는 Newton 방법은 야코비안 $J$로 선형화한 계를 푼다.
\[
J(x_k)s_k=-F(x_k),\qquad x_{k+1}=x_k+s_k
\]
\cref{thm:newton-quadratic}과 같은 논증으로, $J(x_\ast)$가 정칙이고 $J$가 근 근처에서 Lipschitz 연속이면 국소적으로 2차 수렴한다\cite{ortega2000}. 걸음마다 야코비안을 계산하고($n$개의 열을 차분으로 구하면 함수를 $n$번 더 부른다) $O(n^3)$의 LU 분해를 해야 한다. $F(x,y)=(x^2+y^2-4,\ e^x+y-1)$을 $(1,-1.7)$에서 풀면 오차가 $2.3\times10^{-4}$, $1.3\times10^{-8}$이고 세 번째 걸음에서 정확한 근에 닿는다.

\subsection{Broyden's Update}

야코비안을 매번 다시 계산하는 대신, 걸음마다 얻는 정보로 근사 야코비안 $B_k$를 고쳐 쓸 수 있다. 걸음 $s_k=x_{k+1}-x_k$와 함숫값의 변화 $y_k=F(x_{k+1})-F(x_k)$가 있으면 새 근사는 \term{할선 조건}{secant condition} $B_{k+1}s_k=y_k$를 만족해야 한다. 이 조건만으로는 $B_{k+1}$이 정해지지 않으므로 Broyden은 $B_k$에서 가장 적게 바꾼 것을 골랐다\cite{broyden1965}.

\begin{proposition}
\label{prop:broyden}
할선 조건을 만족하는 행렬 가운데 $\|B-B_k\|_F$를 최소로 하는 것은
\[
B_{k+1}=B_k+\frac{(y_k-B_ks_k)s_k^{\mathsf T}}{s_k^{\mathsf T}s_k}
\]
이다.
\end{proposition}

\begin{proof}
$B_{k+1}s_k=y_k$는 바로 확인된다. $Bs_k=y_k$인 임의의 $B$에 대해 $y_k-B_ks_k=(B-B_k)s_k$이므로
\[
B_{k+1}-B_k=(B-B_k)\frac{s_ks_k^{\mathsf T}}{s_k^{\mathsf T}s_k}
\]
이다. 뒤의 인자는 $s_k$ 방향으로의 직교 사영이어서 2-노름이 $1$이고, $\|XP\|_F\le\|X\|_F\|P\|_2$이므로 $\|B_{k+1}-B_k\|_F\le\|B-B_k\|_F$이다.
\end{proof}

갱신은 랭크 $1$이므로, $B_k$의 LU나 QR 인자를 $O(n^2)$에 고칠 수 있다. Broyden 방법은 초선형으로 수렴한다\cite{dennis1977quasi}. 위의 예에서 첫 야코비안만 정확히 주고 그 뒤로 Broyden 갱신을 쓰면 오차가 $2.3\times10^{-4}$, $1.7\times10^{-6}$, $1.0\times10^{-9}$, $3.9\times10^{-12}$, $8.0\times10^{-16}$으로 Newton 방법보다 두 걸음 늦게 닿는다. 대신 야코비안은 한 번만 계산했다.

\subsection{Powell's Hybrid Method}

Newton 방법과 Broyden 방법은 국소적이다. Powell은 $\|F\|^2/2$을 줄이는 방향으로 전역적인 안전장치를 더했다\cite{powell1970hybrid}. 걸음마다 신뢰 영역의 반지름 $\Delta$ 안에서, Newton 걸음이 들어가면 그것을 쓰고 아니면 Newton 걸음과 가장 가파른 내리막 걸음을 잇는 꺾은선(dogleg)이 반지름과 만나는 점을 쓴다. 실제로 줄어든 $\|F\|$와 선형 모형이 예측한 감소를 견주어 반지름을 늘리거나 줄인다. 야코비안은 Broyden 갱신으로 고치다가 진전이 없으면 다시 계산한다. 신뢰 영역과 dogleg은 45장에서 자세히 본다. Moré, Garbow, Hillstrom의 MINPACK이 이 방법을 \code{hybrd}와 \code{hybrj}로 구현했고\cite{more1980minpack} scipy의 \code{root}의 기본 방법이 그것이다.

\section{Root Finding in stardust}
\label{sec:roots-stardust}

\begin{instardust}[\code{docs/design/optimize.md} §3.2, §7]
stardust의 \code{root\_scalar}는 scipy의 Brent 방법(\code{brentq}), TOMS 748, 이분법, Newton, Halley, 할선법을 계산 순서까지 옮겼다. 설계 문서에 따르면 scipy의 C 코드는 clang의 기본 설정으로 빌드되어 $ab+c$ 꼴을 한 번의 반올림(FMA, 8장)으로 계산한다. stardust는 전역 축약을 끄므로 그 자리를 \code{std::fma}로 명시해야 했다. 외삽의 분자 하나를 \code{fma}로 적고서야 $204$개의 시험 문제 모두에서 함수를 부르는 점이 scipy와 비트까지 같아졌다. FMA 없이 계산하면 $12$문제의 호출점이 끝자리에서 갈리고 $7$문제의 근이 달랐다.
\end{instardust}

\begin{instardust}[\code{docs/design/optimize.md} §7]
근의 참값은 \cref{sec:bracket}의 binary64 이분법으로 만든다. Arb 구간 산술로 $f$의 부호를 확정하며 비트 순서로 이분해 각 문제의 근을 이웃한 두 binary64 수 사이로 좁힌다. 이 참값으로 재면 허용오차를 $0$으로 둔 Brent 방법은 $197$문제에서 근을 끼는 두 수에서 $3$ulp 안, TOMS 748은 $199$문제에서 $1$ulp 안에 수렴한다. 방정식계의 \code{root(method='hybr')}는 MINPACK의 산술을 그대로 옮겨 $55$개의 계에서 scipy와 함수를 부른 점이 차례대로 비트까지 같다. 다만 시험점에서 $F$가 NaN이면 scipy는 그 값으로 Broyden 갱신을 해 QR 인자 전체가 NaN이 된다. stardust는 그 갱신을 건너뛴다.
\end{instardust}

\section{Concluding Remarks}

근 찾기에서는 확실함과 빠름을 함께 얻기 어렵다. 이분법은 걸음 수가 미리 정해질 만큼 확실해도 1차 수렴에 그친다. Newton 방법과 할선법은 빠른 대신 시작점이 나쁘면 어디로 갈지 모른다. Brent의 방법과 TOMS 748은 근을 낀 구간을 지키면서 보간 걸음의 속도를 얻어 실제로 쓰이는 1차원 근 찾기의 표준이 되었다. 어떤 방법이든 도달할 수 있는 정확도는 근의 조건수와 함수를 계산하는 방식이 정한다. 다음 장에서는 함수의 근 대신 적분을 계산한다. 거기서도 다항식 보간과 Chebyshev 점이 다시 중요한 구실을 한다.

한 변수 근 찾기의 고전은 Brent의 책\cite{brent1973}이고, 방정식계는 Dennis와 Schnabel\cite{dennis1996}, Ortega와 Rheinboldt\cite{ortega2000}의 책이 표준이다.

\chapterreferences
